% Source: ACM_Templates.pdf pp. 40-66. User eval/pushdown convention; affine composition derived explicitly.
\subsection{线段树：区间乘加与取模}
\label{ds:seg-affine}
\AlgorithmInfo
    {维护 $a_i\gets a_i\cdot mul+add\pmod P$，统一加、乘和赋值。}
    {建树 $O(n)$；修改、查询 $O(\log n)$；空间 $O(n)$。}
    {$n\ge1$，$1\le P\le2^{63}-1$，不要求模数为素数；使用 GNU \InlineCode{__int128}。}
\subsubsection*{原理}
懒标记表示仿射变换 $f(x)=mx+a$。已有标记为 $(m,a)$，随后再施加 $(M,A)$，组合为
\[
 x\longmapsto (xm+a)M+A=x(mM)+(aM+A).
\]
区间和变为 $sum\cdot M+len\cdot A$，标记变为 $(mM,aM+A)$，计算都取模。标记顺序不能交换；单位标记为 $(1,0)$，赋值等价于 $(0,v)$。

\subsubsection*{使用教程}
构造 \InlineCode{AffineSegTree(n,P)}，填 \InlineCode{w[1..n]} 后调用 \InlineCode{build(1,1,n)}；\InlineCode{modify(1,L,R,mul,add)} 对闭区间施加变换，参数可为负数，会归一化。\InlineCode{query(1,L,R).sum} 返回模 $P$ 的区间和；模数不可在建树后直接改动。
\begin{lstlisting}[style=mtycpp,numbers=none]
AffineSegTree seg(4,1000000007);
seg.w={0,1,2,3,4};
seg.build(1,1,4);
seg.modify(1,1,3,1,5); // add 5
seg.modify(1,2,4,3,0); // multiply by 3
seg.modify(1,3,3,0,7); // assign 7
long long ans=seg.query(1,1,4).sum;
\end{lstlisting}
\needspace{7\baselineskip}
\subsubsection*{模板代码}
\CodeFile{code/data-structures/seg_affine.cpp}
\subsubsection*{需要注意}
乘法、合并加法用 128 位中间值，再取模；$P=1$ 也按同一接口处理。
